\begin{figure}
  \centering
  \includegraphics[width=\linewidth]{tornado.pdf}
  \caption{Cost change needed to move the optimal schedule (100~MW new solar, minimising production cost). Each bar spans the changes in that technology's production cost over which the baseline schedule remains optimal; its ends are the breakevens, annotated as percentage changes, beyond which a different schedule is strictly cheaper. The axis is logarithmic, so that a fall and a rise by the same factor are equally far from the centre. A bar running to the plot edge, marked ``no flip'', means the schedule does not change in that direction out to a factor of 10. Technologies are ordered by their nearer breakeven, most fragile first. The breakevens differ somewhat between solar scenarios; see the supplementary table.}
  \label{fig:tornado}
\end{figure}

\begin{figure}
  \centering
  \includegraphics[width=\linewidth]{flip_consequences.pdf}
  \caption{What the first flip is worth. For each technology and each direction, the change in total emissions once the schedule has flipped at the breakeven of Figure~\ref{fig:tornado}, evaluated at true cost coefficients, against the amount of new solar capacity built. Left: the technology becomes cheaper. Right: it becomes dearer. Absent series have no breakeven in that direction. Read together with Figure~\ref{fig:tornado}: a breakeven that is close by is not necessarily one that matters.}
  \label{fig:consequences}
\end{figure}
